\documentclass[5pt, a4paper]{article}

\usepackage[utf8]{inputenc}
\usepackage[T1]{fontenc}
\usepackage{textcomp}
\usepackage{amsmath, amssymb}
\usepackage{multicol}
\usepackage{proof}
\usepackage{enumitem}
\usepackage{fancyhdr}
\usepackage[margin=0.5in]{geometry}
% figure support
\usepackage{import}
\usepackage{xifthen}
\pdfminorversion=7

\usepackage[most]{tcolorbox}
\usepackage{pdfpages}
\usepackage{transparent}
\newcommand{\incfig}[1]{%
    \def\svgwidth{\columnwidth}
    \import{./figures/}{#1.pdf_tex}
}

\tcbset{
colback=gray!5,
colframe=blue!50!white,
coltitle=black,
arc=0pt,
boxrule=-0.1mm,
fonttitle=\bfseries,
coltitle=black,
left=2.5mm,
leftrule=1mm,
right=3.5mm,
colbacktitle=black!10!white,
toptitle=0.75mm,
bottomtitle=0.5mm,
}
\begin{document}
\begin{multicols}{2}
    \begin{tcolorbox}[title=Accuracy vs Precision]
        \begin{enumerate}
            \item Precision: Number of significant bits stored behind the
                decimal.
            \item Accuracy: How close is our output to the true value. 
        \end{enumerate}
    \end{tcolorbox}
    \begin{tcolorbox}[title= Common Algorithm Time Complexity]
        \begin{enumerate}
            \item Multiplying a matrix with a scalar: $O(n)$
            \item Multiplying a $m \times n$ matrix with a $n \times k$ matrix:
                $O(mnk)$
            \item Dot Product of two vectors of size $n$: $O(n)$ 
            \item Multiplication of tridiagonal square matrices: $O(n)$
            \item Multiplication of upper triangular matrices: $O(n^3)$
            \item Forward Elimination: $O(n^3)$
            \item Backward Subtitution: $O(n^2)$ 
            \item Gramm-Schmidtt: $O(n^3)$ 
            \item Householder: $O(\frac{4}{3}n^3)$
             \item Givens: $O(n^3)$ \\ Better for matrices with special
                 structure.
        \end{enumerate}
    \end{tcolorbox}
    \subsection*{Forward Subtitution and Backward Eliminiation}
    The algorithm for solving simple square matrices. Forwards subtitution
    reduces the matrix into an upper triangular form, and the Backward
    elimination eliminates the non-pivot entries.
    \begin{enumerate}
        \item Forward Subtitution 
            \begin{enumerate}
                \item Iterate over each columns.
                \item For a column $i$, subtract each row below the $i$-th row by
                    the $i$-th row such that the $i$-th row entry on the $i$-th
                    column is zero.
            \end{enumerate}
        \item Backward Elimination
            \begin{enumerate}
                \item Iterate over columns from n to $1$. (decreasing)
                \item For each row above the $i$-th row, subtract it such that
                    there is only one entry in the column.
            \end{enumerate}
    \end{enumerate}

\end{multicols}     
\end{document}
